\section{Law, Discipline, and the United States Calendar}\label{sec:law-calendar}

Because Advent prepares for Christmas with quiet restraint, Catholics often wonder what the Church actually binds under law during December and what she leaves to voluntary devotion. Clarity here frees the conscience from both carelessness and scrupulosity.

\subsection{Fasting, Abstinence, and Penance in Current Law}

Every member of the Christian faithful is bound by divine law to do penance (\textit{Code of Canon Law}, 1983, c.~1249). To unite the Church in common observance, canon~1250 designates two universal categories of penitential time: every Friday of the whole year and the season of Lent. Under canon~1251, abstinence from meat is observed on all Fridays (unless a solemnity falls on a Friday), while both fasting and abstinence are observed on Ash Wednesday and Good Friday. Canon~1252 binds those who have completed their fourteenth year to the law of abstinence, and all adults from the beginning of their eighteenth year until the beginning of their sixtieth year to the law of fasting, while canon~1253 empowers episcopal conferences to determine further norms of fasting and abstinence or substitute other forms of penance, especially works of charity and exercises of piety.

Neither Pope St.~Paul~VI's apostolic constitution \textit{Paenitemini} (17 February 1966, Norms~II, IV, and~VI.1.b) nor canons~1250--1253 of the 1983 Code name Advent as a season of mandatory fasting. In the United States, the National Conference of Catholic Bishops exercised its faculty under \textit{Paenitemini} in its \textit{Pastoral Statement on Penance and Abstinence} (18 November 1966, nos.~7, 17, 22--25) and later confirmed its complementary norm on canons~1252--1253 (approved 1983, amended 1984):

\begin{itemize}
  \item \textbf{Ash Wednesday and Good Friday} remain the two days of universal fasting and abstinence in the United States (fasting binds ages 18 through 59; abstinence binds from age 14 upward).
  \item \textbf{Fridays of Lent} are obligatory days of abstinence from meat.
  \item \textbf{Fridays outside Lent---including the Fridays of Advent---}remain penitential days across the entire year. While the U.S. bishops lifted the canonical obligation of meatless Fridays outside Lent, they explicitly exhorted the faithful to preserve Friday abstinence voluntarily as a first-choice mark of Catholic identity, or to substitute another genuine act of penance, prayer, or charity on every Friday.
  \item \textbf{No universal Advent fast days} bind under current law. By contrast, older American Catholic discipline---recorded in the Third Plenary Council of Baltimore's \textit{Manual of Prayers} (1889) and later governed by canon~1252 of the 1917 Code---prescribed fasting on the Winter Ember Days and the Vigil of Christmas (and in many dioceses on the Fridays of Advent).
\end{itemize}

The absence of statutory Advent fasting does not make Advent a season of indulgence. The Church trusts her children to choose a sober, deliberate self-denial---such as keeping the Fridays of Advent meatless, simplifying meals on Wednesdays and Fridays, fasting from trivial entertainment, or giving the money saved from seasonal luxuries directly to the poor.

\subsection{Confession Before Christmas}

No voluntary discipline prepares the soul for the Nativity more deeply than a candid, unhurried sacramental Confession. St.~Francis de Sales urges Philothea in his \textit{Introduction to the Devout Life} (Part~II, ch.~19) not to postpone confession or hide behind vague generalities, but to approach the sacrament frequently, humbly, and plainly so that grace may cleanse the soul's wounds before they harden. Going to confession in the Second or Third Week of Advent---rather than waiting for crowded lines on December~23---allows the last days before Christmas to unfold in peace. Appendix~B prints the classic Act of Contrition from the 1889 \textit{Manual of Prayers} for use in nightly examination of conscience and before entering the confessional.

\subsection{Liturgical Restraint in the Postconciliar Missal}

Even before a homily is preached, the postconciliar Roman Missal teaches the spirit of Advent through four visible and audible signs:

\begin{enumerate}
  \item \textbf{The omission of the \textit{Gloria}.} Under the \textit{General Instruction of the Roman Missal} (GIRM, no.~53), the \textit{Gloria in excelsis Deo} is not sung or said on the Sundays or weekdays of Advent. The Church holds back the angels' Bethlehem hymn (Luke~2:14) so that it breaks forth afresh on Christmas Night. Two Advent occasions interrupt this silence in the United States: the \textit{Gloria} is sung on December~8 (Solemnity of the Immaculate Conception) and on December~12 (Feast of Our Lady of Guadalupe), as well as on November~30 (Feast of St.~Andrew) whenever it falls on a weekday in Week~I (when November~30 coincides with the First Sunday of Advent, the Sunday prevails and the Feast is omitted).
  \item \textbf{The retention of the \textit{Alleluia}.} Unlike Lent, Advent never drops the \textit{Alleluia} before the Gospel (GIRM, no.~62). That single rubric reveals the soul of the season: Advent is not a mourning fast for sin alone, but a watchful, hopeful longing already lit from within by the resurrection of Christ.
  \item \textbf{Violet and rose vestments.} Violet is the proper liturgical color of Advent (GIRM, no.~346), expressing sober watchfulness and preparation. On the Third Sunday of Advent (\textit{Gaudete} Sunday), rose vestments may be worn where it is the practice, lightening the violet at the season's midpoint as the Nativity draws near.
  \item \textbf{Moderation in flowers and music.} Floral decoration of the altar is kept moderate throughout Advent in harmony with the character of the season, avoiding any anticipation of the full festal joy of the Nativity of the Lord (GIRM, no.~305). Likewise, the organ and other musical instruments are played with measured restraint suited to the season's expectant character (GIRM, no.~313).
\end{enumerate}

From December~17 through December~24, the calendar enforces an additional safeguard: these eight privileged weekdays stand at no.~9 in the Table of Liturgical Days (\textit{Universal Norms}, no.~59), outranking every memorial. If an optional memorial occurs (such as St.~Peter Canisius on December~21 or St.~John of Kanty on December~23), the weekday Mass of Advent must be celebrated, though the priest may take the Collect of the memorial in place of the weekday Collect while retaining the Advent readings and remaining prayers (\textit{General Instruction}, no.~355a).

\subsection{The United States Calendar: December 8 and December 12}

Two great Marian celebrations shape Advent in the dioceses of the United States, and their calendar rules deserve special care.

\subsubsection{December 8: The Immaculate Conception of the Blessed Virgin Mary}

Under canon~1246 \S1 of the 1983 Code, the Solemnity of the Immaculate Conception is one of the ten universal holy days of obligation. In the United States, where Mary under the title of her Immaculate Conception is the nation's principal Patroness, December~8 is always observed as a Holy Day of Obligation---even when December~8 falls on a Saturday or a Monday (unlike January~1, August~15, and November~1, whose Mass obligation is abrogated on Saturdays and Mondays under the United States bishops' 1991--1992 complementary norm).

A special calendar collision arises whenever December~8 falls on a Sunday---the Second Sunday of Advent---which occurs fourteen times between 2020 and 2120 (including 2024, 2030, 2041, and 2047). Because Sundays of Advent hold rank~2 in the Table of Liturgical Days while solemnities of the General Calendar hold rank~3 (\textit{Universal Norms}, no.~59), the Second Sunday of Advent always prevails on December~8, and the Solemnity of the Immaculate Conception is transferred liturgically to Monday, December~9 (\textit{Universal Norms}, no.~5).

For decades, United States liturgical calendars stated that when the Solemnity moved to Monday, December~9, the obligation to attend Mass did not transfer with it. In September~2024, however, the Dicastery for Legislative Texts responded to an inquiry from the chair of the USCCB Committee on Canonical Affairs and Church Governance by stating that the precept transferred to Monday, December~9, 2024. Following further consultation between both dicasteries, the Dicastery for Divine Worship and the Discipline of the Sacraments issued a clarifying Note on 23 January 2025 (Prot.~n.~51/25) establishing the rule going forward: in the event of an occasional transfer of a holy day of obligation due to a collision with a privileged Sunday, \textit{the obligation to attend Mass is not transferred to the day} \textit{ad quem} (the day to which the feast is moved). Thus, when December~8 falls on a weekday or Saturday (as on Tuesday, December~8, 2026), it is a Holy Day of Obligation in the United States; when December~8 falls on the Second Sunday of Advent (next in 2030), the Solemnity is celebrated liturgically on Monday, December~9 without a transferred canonical obligation---though faithful Catholics should always verify their own diocesan notice in any transfer year.

\subsubsection{December 12: Feast of Our Lady of Guadalupe}

On December~12, the dioceses of the United States celebrate Our Lady of Guadalupe, Patroness of the Americas, with the rank of a \textbf{Feast} (Table of Liturgical Days, no.~8), a higher dignity than its Optional Memorial rank on the General Roman Calendar. Because it is a Feast, Mass includes the \textit{Gloria} and proper readings honoring Mary's maternal visitation to St.~Juan Diego at Tepeyac in 1531; it is not a Holy Day of Obligation. When December~12 falls on the Third Sunday of Advent---which happens fifteen times between 2020 and 2120 (including 2021, 2027, 2032, and 2038)---the Sunday prevails and the Feast of Our Lady of Guadalupe is omitted for that liturgical year rather than transferred (\textit{Universal Norms}, no.~60).

\subsubsection{The Saints of Advent}

Alongside December~8 and December~12, the universal and United States calendars scatter luminous companions along the Advent road: St.~Andrew the Apostle (November~30, Feast when it falls strictly after the First Sunday on Monday--Wednesday of Week~I; when November~30 falls on the First Sunday of Advent, as in 2025, 2031, and 2036, the Sunday takes precedence under \textit{Universal Norms} no.~59 and the Feast is omitted), St.~Francis Xavier (December~3, Memorial), St.~John Damascene (December~4, Optional Memorial), St.~Nicholas (December~6, Optional Memorial), St.~Ambrose of Milan (December~7, Memorial), St.~Juan Diego Cuauhtlatoatzin (December~9, Optional Memorial, omitted when Immaculate Conception transfers to December~9), Our Lady of Loreto (December~10, Optional Memorial, added to the General Calendar in 2019), St.~Damasus~I (December~11, Optional Memorial), St.~Lucy (December~13, Memorial), and St.~John of the Cross (December~14, Memorial), followed during the privileged week by the commemorations of St.~Peter Canisius (December~21) and St.~John of Kanty (December~23). None of these saints competes with Advent; each shows a distinct face of watching for the Lord.
